\chapter{Experiments and results}
Another problem to consider is fully held-out windows.
\begin{itemize}
    \item Lower error for smaller ts between linear and exponential because of higher discriminative power
\end{itemize}

1. Phase 1, 2048/400 replication
2. Phase 2, expansion to edm 512/100 performance and comparison with baseline models
3. Further expansion and evaluation of EDM models (against baseline possibly)

\section{Introduction}
In this chapter we will analyse the experiments conducted and the results yielded. Before proceeding it is important to underline the two core questions we are trying to answer here:
\begin{itemize}
    \item Are we able to replicate the results obtained by Kim et al. \cite{kim2025diffusion} and to which extent?
    \item Does the embedding into EDM and the introduction of conditional information improved models' performance and ability to replicate true FTS distribution?
\end{itemize}
\section{Phase 1}
For phase 1 the settings are available at Tab.\ref{tab:training_setup}. On top of these, models were trained for 1000 epochs, similarly to Kim et al. \cite{kim2025diffusion} with batch size of 64. The sampling technique is probability flow ODE, which empirically have been show produing more reliable and replicable results.


[INCLUDE IMAGES THAT SHOW FTS AND RELEVANT OTHER GRAPHS OR TABLE]

General structure:
* Aggregate statistics
* distributional information
* Stylized facts + alpha
* path illustration

For conditional add also:
* CRPS
* Interval Coverage

\section{Phase 2}
In this section it is conducted a comparative study of EDM alongside previous implementation. Specifically, we evaluate the conditioning. In order to have a fair assessment, the models analysed in Phase 1 were re-trained with \(L=512\) and \(stride=100\) and the same Heun's ODE sampler was used.

[Remember of the different window-training instantializationhttps://www.trekkinglecco.com/sentiero-viandante/]
